% Avsnitt 1.7, ur lectures/L01/README.md: målen, kontrollfrågorna och blicken framåt.

\section{Sammanfattning}
\label{c1:sec:review}

\needspace{6\baselineskip}
\subsection*{Det här ska du kunna nu}
\begin{itemize}
  \item Härleda en summa-av-produkter ur en sanningstabell, och sedan bygga och simulera nätet i
    CircuitVerse.
  \item Tillämpa De Morgan för att flytta en inversion genom en grind, och använda det för att
    realisera en funktion i en grinduppsättning som saknar operatorn du utgick från.
  \item Läsa och skriva en kombinatorisk VHDL-modul, en \code{entity} med
    \code{std_logic}-portar och en \code{architecture} som driver dess utgång, och säga varför
    VHDL skiljer de två åt.
  \item Förklara varför en konkurrent tilldelning beskriver en ledning snarare än en sats som körs
    en gång, och varför en signal som ligger kvar på \code{'U'} eller \code{'X'} är ett fel och
    inte ett värde.
  \item Köra en utdelad testbänk och avgöra om din modul klarade den.
\end{itemize}

\needspace{6\baselineskip}
\subsection*{Frågor att testa dig själv med}
\begin{enumerate}
  \item Vad är skillnaden mellan en \code{entity} och en \code{architecture}, och varför skiljer
    VHDL på dem?
  \item Varför är \code{x <= a or b;} inte ”en sats som körs en gång”? Vad beskriver den i
    stället?
  \item I bromsassistenten läggs förarens pedal in med OR sist, i stället för att gå genom
    felhanteringslogiken. Vad går sönder om du kopplar den tvärtom?
\end{enumerate}

\needspace{6\baselineskip}
\subsection*{Blicken framåt}
\Chapref{c2:ch} går vidare med:
\begin{itemize}
  \item Karnaughdiagram: att hitta det \emph{minsta} grindnätet, inte bara ett korrekt.
  \item \code{std_logic_vector}, \code{process} och \code{case}.
  \item Submoduler, demonstrerade på en tvåsiffrig hexadecimal display driven av åtta
    strömbrytare.
  \item Självkontrollerande testbänkar: att verifiera din egen VHDL utan kort.
\end{itemize}
